\documentclass[11pt]{article}
\usepackage{amsmath,amssymb,amsthm}
\usepackage[margin=1in]{geometry}
\usepackage{hyperref}

\newtheorem{proposition}{Proposition}
\newtheorem{remark}[proposition]{Remark}

\title{Dividing polynomials by setting the divisor to one\\
\large An iterative algorithm and its application to bifurcation polynomials}
\author{}
\date{\today}

\begin{document}
\maketitle

\begin{abstract}
Given polynomials $Q(z)$ and $p(z)$ with $p\mid Q$, we want the quotient
$B=Q/p$ without performing multivariate polynomial long division.  The
observation is simple: in the ring $\mathbb{Z}[z]/(p-1)$, the divisor $p$
acts as $1$, so $Q\equiv B\cdot p\equiv B$.  When $\deg B<\deg p$, the
reduction modulo $p-1$ recovers $B$ outright.  When $\deg B\ge\deg p$, an
iterative ``base-$(p{-}1)$'' expansion reconstructs $B$ from a sequence of
low-degree residues, each obtained by one cheap reduction.  We work out the
algorithm, verify it on the period-doubling bifurcation polynomials of the
quadratic map $g_c(z)=z^2+c$, and discuss where it helps and where it does
not.
\end{abstract}

\tableofcontents

%======================================================================
\section{The idea}
%======================================================================

Suppose $Q=B\cdot p$ in $\mathbb{Z}[z,c]$ and we want $B$.  Consider the
ring $R=\mathbb{Z}[z,c]/(p-1)$, where the relation $p(z,c)=1$ holds.  In
this ring,
\[
Q = B\cdot p = B\cdot 1 = B.
\]
So reducing $Q$ modulo $p-1$ gives $B \bmod (p-1)$: a polynomial of degree
$<\deg_z(p)$ in $z$ that agrees with $B$ in $R$.  If $\deg_z B<\deg_z p$,
no information is lost and the result \emph{is} $B$.

\paragraph{A tiny example.}
$Q=z^2-4$, $p=z+2$, so $B=z-2$.  The relation $p=1$ gives $z=-1$.
Reducing: $Q(-1)=1-4=-3$.  Since $B$ is linear, $B(z)=z+b$ with leading
coefficient $1$ (from $Q$ and $p$), and $B(-1)=-3$ forces $b=-2$.  So
$B=z-2$.

%======================================================================
\section{When $\deg B\ge\deg p$: the iterative expansion}
\label{sec:iteration}
%======================================================================

Write $d=\deg_z(p-1)=\deg_z p$.  If $\deg_z B\ge d$, the reduction
$Q\bmod(p{-}1)$ collapses $B$ past what we want.  The fix is a
``base-$(p{-}1)$'' expansion: express $B$ in powers of $(p-1)$,
\begin{equation}\label{eq:expansion}
B = R_0 + R_1\,(p{-}1) + R_2\,(p{-}1)^2 + \cdots + R_N\,(p{-}1)^N,
\end{equation}
where each $R_k$ has $\deg_z R_k < d$ and
$N=\lceil\deg_z B/d\rceil$.  This is analogous to writing an integer in a
given base: each ``digit'' $R_k$ lives in the range $[0,\,(p{-}1))$ (as a
polynomial in $z$), and the expansion is unique.

\subsection{The recursion}

Set $Q_0=Q$ (which equals $B\cdot p$).  Then iterate:
\begin{enumerate}
\item \textbf{Reduce:} $R_k = Q_k \bmod (p{-}1)$.  Since $p\equiv 1$, this
      equals $A_k\bmod(p{-}1)$ where $Q_k=A_k\cdot p$.
\item \textbf{Peel and divide:}
      $Q_{k+1} = \dfrac{Q_k - R_k\cdot p}{p-1}$.
\end{enumerate}
The division in step~2 is exact: $Q_k\equiv R_k\cdot p\pmod{p{-}1}$ (because
$p\equiv 1$ and $Q_k\equiv R_k$), so $(p{-}1)\mid(Q_k-R_k\cdot p)$.  Each
step lowers $\deg_z Q_k$ by $d$, so after $N+1$ residues the process
terminates with $Q_{N+1}=0$.

\subsection{Why it is exact}

The expansion~\eqref{eq:expansion} is the unique representation of $B$ in
the ``$(p{-}1)$-adic'' basis.  Each $R_k$ is the $k$-th digit.  The
recursion extracts digits from the top: $R_0$ is the lowest digit of $B$
(its residue mod~$p{-}1$), and after peeling it off and dividing, $R_1$ is
the next digit, and so on.  Since $B$ has finite degree, the process
terminates.

%======================================================================
\section{Worked examples from the quadratic map}
\label{sec:examples}
%======================================================================

The quadratic map $g_c(z)=z^2+c$ generates iterates $g_c^m(z)$ of degree
$2^m$ in $z$.  The period-$m$ cycle polynomial is $g_c^m(z)-z$, and the
``new period-$m$'' polynomial (excluding lower-period points) is
\[
B = \frac{g_c^m(z)-z}{g_c^{m/2}(z)-z} = \frac{Q}{p}.
\]
Both $Q$ and $p$ are \emph{monic} in $z$, so $p-1$ is also monic and the
reduction rule $z^d=(\text{lower terms})$ involves no denominators.

\subsection{The $3\to6$ bifurcation: $\deg B=6$, $\deg p=2$}

Here $p=g(z)-z=z^2+c-z$, $Q=g^3(z)-z$ (degree~8), $B=Q/p$ (degree~6).
The modulus is $p-1=z^2+c-z-1$, giving the rule $z^2=z-c+1$.  Four
residues are needed ($N=3$):
\begin{center}
\renewcommand{\arraystretch}{1.2}
\begin{tabular}{cl}
$k$ & $R_k$ \\\hline
0 & $13-8c+20z$ \\
1 & $20-4c+18z$ \\
2 & $9+4z$ \\
3 & $1$
\end{tabular}
\end{center}
The residues have $\deg_z R_k<2$ and small $c$-degree (at most~1).  The
reconstruction $B=R_0+R_1(p{-}1)+R_2(p{-}1)^2+R_3(p{-}1)^3$ matches the
ground truth (from FORM's \texttt{div\_}) exactly, with zero remainder at
every step.

\subsection{The $4\to8$ bifurcation: $\deg B=12$, $\deg p=4$}

Here $p=g^2(z)-z$ (degree~4), $Q=g^4(z)-z$ (degree~16), $B$ has degree~12
in $z$ and degree~6 in $c$.  The modulus $p-1$ has degree~4, so
$N=\lceil12/4\rceil=3$ and we need four residues:
\begin{center}
\renewcommand{\arraystretch}{1.2}
\begin{tabular}{ccc}
$k$ & $\deg_z R_k$ & $\deg_c R_k$ \\\hline
0 & 3 & 1 \\
1 & 2 & 1 \\
2 & 1 & 0 \\
3 & 0 & 0
\end{tabular}
\end{center}
Again: exact reconstruction, zero remainder, and the $c$-degree of the
residues stays low (at most~1, despite $B$ having $c$-degree~6).  The
intermediate polynomials $Q_k$ decrease in $z$-degree by~4 each step
(16, 12, 8, 4) and the division by $p-1$ (degree~4) is cheap.

%======================================================================
\section{Cost analysis}
\label{sec:cost}
%======================================================================

The standard division $B=Q/p$ via FORM's \texttt{div\_} performs
$\deg_z Q-\deg_z p$ steps of multivariate polynomial long division.  Each
step involves the full coefficient ring $\mathbb{Z}[c]$, and the
intermediate expressions can be large.

The iterative algorithm replaces this with $N+1$ cycles of:
\begin{enumerate}
\item One reduction mod~$(p{-}1)$ (repeated application of
      $z^d=\text{rhs}$ via \texttt{\#inside} --- very cheap).
\item One polynomial multiplication $R_k\cdot p$ (the residue times the
      divisor).
\item One exact division by $p-1$ (degree~$d$, via \texttt{div\_}).
\end{enumerate}
The total number of degree-$d$ divisions is $N=\lceil\deg_z B/d\rceil$,
compared to one degree-$\deg_z p$ division on a degree-$\deg_z Q$ dividend
in the standard approach.  Each individual division is cheaper (the
dividend $Q_k$ has degree at most $\deg_z Q-k\cdot d$, and the divisor has
degree $d$), but there are $N$ of them.

\paragraph{When it wins.}
The win is largest when $\deg_z Q\gg\deg_z p$ and the coefficient ring
(polynomials in $c$) is heavy: each step of the standard \texttt{div\_}
manipulates the full coefficient bulk, while the iterative steps work on
progressively smaller $Q_k$.  The reduction mod~$(p{-}1)$ via
\texttt{\#inside} is essentially free (no active-expression overhead), and
the degree-$d$ divisions are fast.

\paragraph{When it does not.}
If $\deg_z B$ is only slightly larger than $\deg_z p$ (say, $N=1$ or $2$),
the overhead of the extra multiplications $R_k\cdot p$ and the $N$ separate
\texttt{div\_} calls may exceed the cost of one direct \texttt{div\_}.  The
algorithm also assumes the divisor $p$ is \emph{monic} in $z$; for
non-monic $p$ (as in the logistic map, where the leading coefficient is a
power of $r$), each reduction step would require dividing coefficients in
the parameter ring, complicating the arithmetic.

%======================================================================
\section{Where it fails}
\label{sec:fails}
%======================================================================

\subsection{Non-monic divisors}

The reduction rule $z^d=\text{rhs}$ comes from solving $p-1=0$ for $z^d$.
If $p$ is monic, $p-1$ is monic, and the rule has integer (polynomial)
coefficients.  If $p$ has leading coefficient $\ell(c)\ne1$, the rule
becomes $z^d = -(\text{lower terms})/\ell(c)$, introducing rational
functions of $c$.  The iterative algorithm still works algebraically (the
divisions are exact), but the intermediate expressions live in
$\mathbb{Q}(c)[z]$ rather than $\mathbb{Z}[c,z]$, which FORM does not
handle natively.

For the quadratic map $g_c(z)=z^2+c$, all iterates are monic in $z$, so
this is not an issue.  For the logistic map $f_r(x)=rx(1{-}x)$, the
iterates have leading coefficient $(-r)^{2^m-1}$ in $x$, and the algorithm
would need modification (e.g., pseudo-division, or clearing denominators at
each step).

\subsection{PrimitivePart interference}

The library procedure \texttt{PolyReduce} calls \texttt{PrimitivePart}
after each reduction, which divides out the GCD of numerical coefficients.
This can flip signs or rescale the residues, breaking the exact
reconstruction (each $R_k$ must be the \emph{true} remainder, not a
primitive associate).  The workaround is to bypass \texttt{PolyReduce}
entirely and perform the reduction with a raw \texttt{\#inside} loop:
\begin{verbatim}
    #inside $R
        id z^d = rhs;
    #endinside
\end{verbatim}
repeated until $\deg_z R<d$.  This is both simpler and faster than the
general-purpose \texttt{PolyReduce} for this application.

\subsection{Coefficient growth}

In the worked examples, the $c$-degree of the residues $R_k$ stayed low
(at most~1).  This is not guaranteed in general: each step multiplies $R_k$
by $p$ (which has $c$-degree~1 for the quadratic map), and the division by
$p-1$ can redistribute $c$-powers across terms.  For very high-order
bifurcations, the intermediate $c$-degrees could grow, increasing the cost
of the degree-$d$ divisions.  Whether this growth is a practical bottleneck
depends on the specific polynomial structure and would need to be measured
empirically at higher orders.

%======================================================================
\section{Summary}
\label{sec:summary}
%======================================================================

\begin{center}
\renewcommand{\arraystretch}{1.3}
\begin{tabular}{lcc}
 & Standard \texttt{div\_} & Iterative $(p{=}1)$ \\\hline
Divisions & 1, degree $\deg p$ on degree $\deg Q$ & $N$, degree $d$ on
  decreasing $Q_k$ \\
Requires monic $p$? & no & yes (or pseudo-division) \\
Extra operations & none & $N$ multiplications $R_k\cdot p$ \\
PrimitivePart safe? & yes & no (use raw \texttt{\#inside}) \\
Best regime & small $\deg Q/\deg p$ & large $\deg Q/\deg p$ with heavy
  coefficients
\end{tabular}
\end{center}

The algorithm is a direct analogue of radix conversion: the ``digits'' $R_k$
are the coefficients of $B$ in the base $(p{-}1)$, and the iteration extracts
them one at a time from the product $Q=B\cdot p$ by exploiting the identity
$p\equiv 1\pmod{p{-}1}$.

\end{document}
